% GENERATED FILE. Edit canonical lessons, then run npm run generate:books.
% canonical-source-hash: fnv1a64:f9580d7947e7be7a
% canonical-lessons: ZH-C254-chusheng, ZH-C254-jiehun, ZH-C254-biye, ZH-C254-hunli, ZH-C254-shenghuo

\chapter{Born and Married}
\label{ch:zh-born-and-married}
\hlchaptermodality{254}

\begin{chapteropening}
\textbf{By the end of this chapter:} \emph{I can say a baby has been born, someone has married or finished school, and talk about a wedding and daily life.}

Complete the last lesson of chapter 254: I can say a baby has been born, someone has married or finished school, and talk about a wedding and daily life.
\end{chapteropening}

\section[\texorpdfstring{\zh{出生} (chūshēng)}{出生 (chūshēng)}]{\zh{出生} (chūshēng) — to be born}

\label{lesson:ZH-C254-chusheng}



\begin{quote}
Before the new one: say the Mandarin for sons and daughters, then the Mandarin for a relative.
\end{quote}

\subsection*{You'll want to know: \zh{出生}}
\textbf{\zh{出生}} — \emph{chūshēng} — \textquotedblleft{}to be born\textquotedblright{}.

\begin{quote}
\textbf{\zh{他的孩子出生了}} \emph{tā de háizi chūshēng le} — His baby has been born.
\end{quote}

\begin{cousinweb}
\textbf{\zh{出生}} (\emph{chūshēng}) is built from \zh{出} 'to come out' and \zh{生} 'to grow', both inherited from Middle Chinese, \emph{tsyhwit} and \emph{srjaeng} (Guangyun, in Baxter's transcription), which glosses the second 'to grow'.
\end{cousinweb}

\begin{grammarlens}[title={What you've built}]
One more word for talking about the past.
\end{grammarlens}

\subsection*{Your turn}
\begin{itemize}
\raggedright
  \item \emph{Say it:} \emph{chūshēng}
  \item \emph{Say it:} \emph{chūshēng}, once more
  \item \emph{Say it:} \emph{tā de háizi chūshēng le} — His baby has been born.
  \item \emph{Recall:} say the Mandarin for sons and daughters, then the Mandarin for a relative, then say \emph{chūshēng} again
\end{itemize}

\begin{tcolorbox}[breakable,colback=teal!4,colframe=teal!35!black,title={Before you move on}]
What does \zh{出生} mean? (\textquotedblleft{}To be born\textquotedblright{}.) Say it once more.
\end{tcolorbox}
\section[jiéhūn]{jiéhūn — to get married}

\label{lesson:ZH-C254-jiehun}



\begin{quote}
Before the new one: say the Mandarin for a relative, then the Mandarin for to be born.
\end{quote}

\subsection*{You'll want to know: jiéhūn}
\emph{jiéhūn} — \textquotedblleft{}to get married\textquotedblright{}.

\begin{quote}
\emph{wǒ gēge jiéhūn le} — My brother has got married.
\end{quote}

Its characters come later; for now it is a word for the ear.

\begin{cousinweb}
\emph{jiéhūn} is built from \emph{jié} 'to tie' and \emph{hūn} 'marriage', both inherited from Middle Chinese, \emph{ket} and \emph{xwon} (Guangyun, in Baxter's transcription), which glosses the first 'to tie' and the second 'marriage'.
\end{cousinweb}

\begin{grammarlens}[title={What you've built}]
One more word for talking about the past.
\end{grammarlens}

\subsection*{Your turn}
\begin{itemize}
\raggedright
  \item \emph{Say it:} \emph{jiéhūn}
  \item \emph{Say it:} \emph{jiéhūn}, once more
  \item \emph{Say it:} \emph{wǒ gēge jiéhūn le} — My brother has got married.
  \item \emph{Recall:} say the Mandarin for a relative, then the Mandarin for to be born, then say \emph{jiéhūn} again
\end{itemize}

\begin{tcolorbox}[breakable,colback=teal!4,colframe=teal!35!black,title={Before you move on}]
What does jiéhūn mean? (\textquotedblleft{}To get married\textquotedblright{}.) Say it once more.
\end{tcolorbox}
\section[bìyè]{bìyè — to graduate; to finish school}

\label{lesson:ZH-C254-biye}



\begin{quote}
Before the new one: say the Mandarin for to be born, then the Mandarin for to get married.
\end{quote}

\subsection*{You'll want to know: bìyè}
\emph{bìyè} — \textquotedblleft{}to graduate; to finish school\textquotedblright{}.

\begin{quote}
\emph{tā míngnián dàxué bìyè} — She finishes university next year.
\end{quote}

Its characters come later; for now it is a word for the ear.

\begin{cousinweb}
\emph{bìyè} is built from \emph{bì} 'to finish' and \emph{yè} 'a task', both inherited from Middle Chinese, \emph{pjit} and \emph{ngjaep} (Guangyun, in Baxter's transcription), which glosses the first 'to end' and the second 'an affair'.
\end{cousinweb}

\begin{grammarlens}[title={What you've built}]
One more word for talking about the past.
\end{grammarlens}

\subsection*{Your turn}
\begin{itemize}
\raggedright
  \item \emph{Say it:} \emph{bìyè}
  \item \emph{Say it:} \emph{bìyè}, once more
  \item \emph{Say it:} \emph{tā míngnián dàxué bìyè} — She finishes university next year.
  \item \emph{Recall:} say the Mandarin for to be born, then the Mandarin for to get married, then say \emph{bìyè} again
\end{itemize}

\begin{tcolorbox}[breakable,colback=teal!4,colframe=teal!35!black,title={Before you move on}]
What does bìyè mean? (\textquotedblleft{}To graduate; to finish school\textquotedblright{}.) Say it once more.
\end{tcolorbox}
\section[hūnlǐ]{hūnlǐ — wedding}

\label{lesson:ZH-C254-hunli}



\begin{quote}
Before the new one: say the Mandarin for to get married, then the Mandarin for to graduate.
\end{quote}

\subsection*{You'll want to know: hūnlǐ}
\emph{hūnlǐ} — \textquotedblleft{}wedding\textquotedblright{}.

\begin{quote}
\emph{wǒmen qù cānjiā tāmen de hūnlǐ} — We're going to their wedding.
\end{quote}

Its characters come later; for now it is a word for the ear.

\begin{cousinweb}
\emph{hūnlǐ} is built from \emph{hūn} 'marriage' and \emph{lǐ} 'a rite', both inherited from Middle Chinese, \emph{xwon} and \emph{lejX} (Guangyun, in Baxter's transcription), which quotes an older dictionary for the second: what is done to serve the spirits.
\end{cousinweb}

\begin{grammarlens}[title={What you've built}]
One more word for talking about the past.
\end{grammarlens}

\subsection*{Your turn}
\begin{itemize}
\raggedright
  \item \emph{Say it:} \emph{hūnlǐ}
  \item \emph{Say it:} \emph{hūnlǐ}, once more
  \item \emph{Say it:} \emph{wǒmen qù cānjiā tāmen de hūnlǐ} — We're going to their wedding.
  \item \emph{Recall:} say the Mandarin for to get married, then the Mandarin for to graduate, then say \emph{hūnlǐ} again
\end{itemize}

\begin{tcolorbox}[breakable,colback=teal!4,colframe=teal!35!black,title={Before you move on}]
What does hūnlǐ mean? (\textquotedblleft{}Wedding\textquotedblright{}.) Say it once more.
\end{tcolorbox}
\section[shēnghuó]{shēnghuó — life; daily life}

\label{lesson:ZH-C254-shenghuo}



\begin{quote}
Before the new one: say the Mandarin for to graduate, then the Mandarin for a wedding.
\end{quote}

\subsection*{You'll want to know: shēnghuó}
\emph{shēnghuó} — \textquotedblleft{}life; daily life\textquotedblright{}.

\begin{quote}
\emph{zhèr de shēnghuó hěn hǎo} — Life here is good.
\end{quote}

Its characters come later; for now it is a word for the ear.

\begin{cousinweb}
\emph{shēnghuó} is built from \zh{生} 'to grow' and \emph{huó} 'alive', both inherited from Middle Chinese, \emph{srjaeng} and \emph{hwat} (Guangyun, in Baxter's transcription), which glosses the second 'not dead'.
\end{cousinweb}

\begin{grammarlens}[title={What you've built}]
One more word for talking about the past.
\end{grammarlens}

\subsection*{Your turn}
\begin{itemize}
\raggedright
  \item \emph{Say it:} \emph{shēnghuó}
  \item \emph{Say it:} \emph{shēnghuó}, once more
  \item \emph{Say it:} \emph{zhèr de shēnghuó hěn hǎo} — Life here is good.
  \item \emph{Recall:} say the Mandarin for to graduate, then the Mandarin for a wedding, then say \emph{shēnghuó} again
\end{itemize}

\begin{tcolorbox}[breakable,colback=teal!4,colframe=teal!35!black,title={Before you move on}]
What does shēnghuó mean? (\textquotedblleft{}Life; daily life\textquotedblright{}.) Say it once more.
\end{tcolorbox}
